\documentclass[aps,prl,twocolumn,showpacs,preprintnumbers]{revtex4-2}

\usepackage{amsmath,amssymb,amsfonts}
\usepackage{xcolor}
\usepackage{booktabs}
\usepackage{microtype}
\usepackage{hyperref}

\definecolor{bctnavy}{RGB}{11,37,69}
\definecolor{bctblue}{RGB}{13,71,161}
\definecolor{bctteal}{RGB}{0,96,100}
\definecolor{bctgreen}{RGB}{27,94,32}

\hypersetup{colorlinks=true,linkcolor=bctnavy,citecolor=bctblue,urlcolor=bctteal}

\newcommand{\roct}{r_{\mathrm{oct}}}
\newcommand{\rtet}{r_{\mathrm{tet}}}
\newcommand{\azero}{\alpha_0}
\newcommand{\lP}{\ell_P}
\newcommand{\dBCT}{d_{\mathrm{BCT}}}

\begin{document}

\preprint{BCT Letter 129 --- 24 March 2026}

\title{Distance Is Emergent:\\
Macroscopic Geometry from BCT Void-Hop Counting}

\author{Michel Robert Cabri\'e}
\affiliation{Independent Researcher \& Artist, Victoria, Australia}
\email{ZeroFreeParameters@gmail.com}
\homepage{https://patreon.com/cw/TheBCTSuperfluidLatticeModel}

\date{24 March 2026}

\begin{abstract}
Distance is not a fundamental property of nature. Within the BCT
Superfluid Lattice Model, macroscopic distance emerges as the
coarse-grained count of void-to-void phonon transitions in the
Planck-scale D4 condensate. The two primitive length scales are
the octahedral and tetrahedral void inradii,
$\roct = (\sqrt{2}-1)/2$ and $\rtet = (\sqrt{6}-2)/4$ (in Planck
units), which are not distances between points but geometric
properties of the condensate structure itself. The Minkowski metric
$ds^2 = c^2dt^2 - d\mathbf{x}^2$ is derived --- not assumed ---
as the long-wavelength limit of the D4 phonon dispersion relation,
which is forced to be isotropic by Lorentz invariance
(the BCT Uniqueness Theorem, Letter~18). BCT predicts a minimum
physically meaningful length $\ell_{\min} = \rtet\,\lP =
0.11237\,\lP = 1.816\times10^{-36}$~m, below which the concept
of distance dissolves into void geometry. The granularity of
distance measurement is set by $\azero = \roct\rtet/\pi$,
the BCT vacuum coupling.
\end{abstract}

\maketitle

% ─────────────────────────────────────────────────────────────────────
\section{The Problem with Distance}
% ─────────────────────────────────────────────────────────────────────

Every framework in fundamental physics assumes distance.
General relativity postulates a metric tensor $g_{\mu\nu}$
from which distances are computed. Quantum field theory assumes
a background Minkowski spacetime with its inherited notion of
separation. String theory postulates a target space.
Even loop quantum gravity, which discretises area and volume,
still assumes the continuous manifold as a starting point.

None of these frameworks \emph{derives} distance.
They assume it and work outward.

The BCT Superfluid Lattice Model takes the opposite approach.
Starting from the single geometric postulate that the quantum
vacuum is a Planck-scale D4 superfluid lattice with axial ratio
$c/a = \sqrt{2}$ (forced by Lorentz invariance, Letter~18~\cite{letter18}),
distance is \emph{derived} as an emergent, counting-based property
of the condensate.

% ─────────────────────────────────────────────────────────────────────
\section{The Primitive BCT Length Scales}
% ─────────────────────────────────────────────────────────────────────

The BCT D4 lattice contains two types of void: octahedral and
tetrahedral. Their inradii in Planck units are~\cite{letter1}:
\begin{align}
  \roct &= \frac{\sqrt{2}-1}{2} = 0.20711,
  \label{eq:roct}\\
  \rtet &= \frac{\sqrt{6}-2}{4} = 0.11237.
  \label{eq:rtet}
\end{align}
These are \emph{not} distances between points. They are geometric
properties of the condensate --- the radii of the largest spheres
that fit inside each void type. They are determined entirely by
the D4 packing geometry and the Planck length; no additional
input is required.

The BCT vacuum coupling $\azero = \roct\rtet/\pi = 0.0074081$
is the product of these two scales, normalised by $\pi$ (the
Josephson phase of the condensate~\cite{letter27}).

% ─────────────────────────────────────────────────────────────────────
\section{How Distance Emerges}
% ─────────────────────────────────────────────────────────────────────

\subsection{Phonons as distance measurers}

A particle propagating through the BCT condensate is a phonon ---
a quantised excitation of the D4 superfluid. It propagates by
exciting successive void sites: each hop carries the excitation
from one void centre to a neighbouring void centre, a displacement
of one lattice constant $a = \lP$ (one Planck length).

\emph{Macroscopic distance is the count of these hops.}

For a phonon traversing $N_{\mathrm{oct}}$ octahedral void sites
and $N_{\mathrm{tet}}$ tetrahedral void sites along a path:
\begin{equation}
  \dBCT = \bigl(N_{\mathrm{oct}}\,\roct + N_{\mathrm{tet}}\,\rtet\bigr)
  \,\lP.
  \label{eq:dBCT}
\end{equation}
This is a discrete, integer-valued counting operation.
It becomes the familiar continuous distance in the limit
$N \to \infty$, where fluctuations between oct and tet hops
average out and the granularity $\azero\lP$ becomes negligible.

\subsection{The metric is derived}

In standard physics, the Minkowski metric
\begin{equation}
  ds^2 = c^2\,dt^2 - dx^2 - dy^2 - dz^2
  \label{eq:minkowski}
\end{equation}
is postulated. In BCT, it is \emph{derived} from the D4 phonon
dispersion relation. The dispersion relation of the BCT condensate
(Letter~18) is:
\begin{equation}
  \omega^2(\mathbf{k}) = c^2\,|\mathbf{k}|^2,
  \label{eq:dispersion}
\end{equation}
which is isotropic --- the same speed $c$ in all directions ---
because the axial ratio $c/a = \sqrt{2}$ is the \emph{unique}
value that gives isotropic phonon propagation in a BCT lattice.
This uniqueness is proven in Letter~18 (the BCT Uniqueness Theorem).

The low-energy effective field theory of an isotropic relativistic
condensate with dispersion~\eqref{eq:dispersion} has the Poincar\'e
group ISO(3,1) as its symmetry group, and the resulting
spacetime metric is precisely Eq.~\eqref{eq:minkowski}.

\textbf{The Minkowski metric is not fundamental. It is the
long-wavelength limit of void-hop counting in the D4 condensate.}

% ─────────────────────────────────────────────────────────────────────
\section{The Minimum Distance}
% ─────────────────────────────────────────────────────────────────────

Below the tetrahedral void inradius $\rtet\lP$, the BCT condensate
has no structure. There are no void sites; there are no hop
transitions; there is no counting operation. The concept of
``distance'' has no referent.

\textbf{BCT Prediction~\#170:}
\begin{equation}
  \ell_{\min} = \rtet\,\lP
  = \frac{\sqrt{6}-2}{4}\,\lP
  = 0.11237\,\lP
  = 1.816\times10^{-36}\;\text{m}.
  \label{eq:lmin}
\end{equation}
This is the smallest length with physical meaning.
Below $\ell_{\min}$, neither distance nor geometry are defined;
only the abstract D4 packing structure exists.

Note that $\ell_{\min}$ is \emph{smaller} than the Planck length
by the factor $\rtet = 0.11237$. This is a specific, parameter-free
BCT prediction distinguishing it from frameworks that identify
$\lP$ itself as the minimum length.

% ─────────────────────────────────────────────────────────────────────
\section{The Granularity of Distance}
% ─────────────────────────────────────────────────────────────────────

Between $\ell_{\min}$ and $\roct\lP$, only tetrahedral hops are
possible. Above $\roct\lP$, both oct and tet hops contribute.
The transition between these regimes occurs at:
\begin{equation}
  \ell_{\mathrm{oct}} = \roct\,\lP
  = \frac{\sqrt{2}-1}{2}\,\lP
  = 0.20711\,\lP
  = 3.347\times10^{-36}\;\text{m}.
  \label{eq:loct}
\end{equation}

The \emph{granularity} of distance measurement --- the quantum
of action associated with a distance measurement --- is set by
$\azero = \roct\rtet/\pi$. A distance measurement carried out
with resolution $\Delta d$ has an associated action
$\Delta S \sim \hbar\,\azero$ when $\Delta d \sim \ell_{\min}$.
This connects the discreteness of distance to the discreteness
of action (Planck's constant) through the BCT vacuum geometry.

% ─────────────────────────────────────────────────────────────────────
\section{Relation to Other Frameworks}
% ─────────────────────────────────────────────────────────────────────

\begin{table}[h]
\centering
\caption{BCT vs other discrete spacetime frameworks.}
\label{tab:compare}
\begin{tabular}{lll}
\toprule
Framework & Minimum length & Free parameters \\
\midrule
Loop Quantum Gravity  & $\sim\lP$ (derived)    & Many \\
Causal Set Theory     & $\sim\lP$ (postulated) & Several \\
Doubly Special Rel.   & $\lP$ (postulated)     & 1 (DSR scale) \\
String Theory         & $\ell_s$ (free)        & Many \\
\textbf{BCT}          & $\rtet\lP$ (derived)   & \textbf{Zero} \\
\bottomrule
\end{tabular}
\end{table}

BCT agrees with loop quantum gravity (LQG) and causal set theory
that spacetime distance is discrete and emergent, not fundamental.
It differs from all of them in deriving the specific value of the
minimum length from the D4 packing geometry, with zero free
parameters.

The BCT picture also accords with the holographic principle:
the void boundaries \emph{are} the holographic screens. Distance
between two regions is encoded in the boundary structure of the
condensate separating them --- exactly as holography requires ---
but with a specific, derivable microscopic realisation.

% ─────────────────────────────────────────────────────────────────────
\section{Why Distance Seems Continuous}
% ─────────────────────────────────────────────────────────────────────

At human scales, distances involve $N \sim 10^{35}$ Planck-length
hops. The ratio $\rtet/\roct = 0.543$ means that the BCT distance
``quantum'' alternates between two values, but on any macroscopic
path the distribution of oct and tet hops averages to a smooth
continuum. The relative fluctuation scales as $1/\sqrt{N} \sim
10^{-17}$, far below any conceivable measurement precision.

Distance \emph{appears} continuous because we are not sensitive
to the granularity. It \emph{is} discrete, emergent, and grounded
in sphere-packing geometry.

% ─────────────────────────────────────────────────────────────────────
\section{The Philosophical Punchline}
% ─────────────────────────────────────────────────────────────────────

The question ``why does distance exist?'' reduces, in BCT, to:
``why does the D4 lattice have two void types?''

The answer: because the D4 lattice is the densest sphere packing
in four dimensions (Viazovska~2016~\cite{viazovska}), and the
densest packing in 4D has exactly this void structure.

And the question ``why is the D4 lattice the ground state?''
reduces to: ``why is the universe 4-dimensional?'' --- which is
the question BCT correctly identifies as pointing beyond physics
into pure mathematics.

Distance exists because sphere-packing in four dimensions has
two types of interstice. That is the deepest answer physics
can give.

\begin{acknowledgments}
This Letter was derived in response to the observation that
distance is not a fundamental concept, made during a session
on 24~March~2026. The BCT programme agrees: distance is not
fundamental. It is the count of void-hop transitions in the
D4 condensate. This was always implied by the programme;
this Letter makes it explicit.
\end{acknowledgments}

\begin{thebibliography}{9}

\bibitem{letter1}
M.~R.~Cabri\'e, \textit{BCT Letter 1: The BCT Vacuum Postulate},
Zenodo (2026). 10.5281/zenodo.18958207

\bibitem{letter18}
M.~R.~Cabri\'e, \textit{BCT Letter 18: The BCT Uniqueness Theorem
--- $c/a=\sqrt{2}$ from Lorentz Invariance},
Zenodo (2026). 10.5281/zenodo.18959915

\bibitem{letter27}
M.~R.~Cabri\'e, \textit{BCT Letter 27: The Gross-Pitaevskii
Axiom as Theorem}, Zenodo (2026).

\bibitem{letter128}
M.~R.~Cabri\'e, \textit{BCT Letter 128: The Cabri\'e Chain of
Necessity}, Zenodo (2026).

\bibitem{viazovska}
M.~S.~Viazovska,
\textit{The sphere packing problem in dimension~8},
Ann.\ Math.\ \textbf{185}, 991 (2017).

\bibitem{rovelli}
C.~Rovelli \& L.~Smolin,
\textit{Discreteness of area and volume in quantum gravity},
Nucl.\ Phys.\ B \textbf{442}, 593 (1995).

\bibitem{bombelli}
L.~Bombelli, J.~Lee, D.~Meyer \& R.~Sorkin,
\textit{Spacetime as a causal set},
Phys.\ Rev.\ Lett.\ \textbf{59}, 521 (1987).

\end{thebibliography}

\end{document}
